\slajdid{09_2-s012}
\begin{frame}[label=09_2-s012]{p-kanalni MOS}
\gusttekst\setlength{\abovedisplayskip}{3pt}\setlength{\belowdisplayskip}{3pt}\setlength{\abovedisplayshortskip}{2pt}\setlength{\belowdisplayshortskip}{2pt}\textbf{Prvi način pisanja – „svi naponi su negativni – označavanje kao kod n-kanalnog, struja ide od sorsa ka drejnu“}
\[I_{Gp}=0\]
\[I_{Dp}=I_{Sp}\ge0\quad\textit{za}\quad V_{GSp}\le V_{Tp}\quad(V_{Tp}<0)\]
\[V_{DSpsat}=\frac{(V_{GSp}-V_{Tp})}{1+\frac{(V_{GSp}-V_{Tp})}{L_pE_{Cp}}}=\frac{(V_{GSp}-V_{Tp})L_pE_{Cp}}{L_pE_{Cp}+(V_{GSp}-V_{Tp})}\quad(E_{Cp}<0)\]
Tranzistor vodi u triodnoj, omskoj, neaktivnoj oblasti, kada je
\[V_{DSp}\ge V_{DSpsat}\]
i tada je njegova struja data izrazom
\[I_{Dp}=\frac{k_p}2\frac1{1+\frac{V_{DSp}}{L_pE_{Cp}}}\left(2V_{DSp}(V_{GSp}-V_{Tp})-V_{DSp}^2\right)\]
\end{frame}
